\documentclass[a4paper,10pt,twocolumn,landscape]{article}
	\usepackage[ansinew]{inputenc}
	\usepackage{ae}
	\usepackage[brazil]{babel}
	\usepackage{xcolor}
	\usepackage{calc}
	\usepackage{wallpaper}
	\usepackage[margin=2cm]{geometry}
	\usepackage{SIunits}
	\usepackage{amsmath}
	\usepackage{paralist}

	% Minhas definições
	\newcounter{exercicio}
\newenvironment{exercicio}{%
	\refstepcounter{exercicio}%
	\penalty-200
	\noindent\colorbox{blue!60!black}{\makebox[\columnwidth-\fboxsep*2][c]{\textbf{\color{white}Exercício~\theexercicio}}}\smallskip
}{\par\medskip}
\newcommand{\foreign}[1]{\textsl{#1}}
\newcommand{\digite}[1]{{\fontfamily{cmss}\fontseries{bx}\selectfont#1}}	
\newcommand{\cs}[1]{{\normalfont\textbackslash\color{blue!50!black}#1}}
\newcommand{\pkg}[1]{{\normalfont\sffamily\color{orange}#1}}
\newcommand{\env}[1]{{\normalfont\sffamily\color{green!50!black}#1}}
\let\comando=\cs
\let\package=\pkg
\let\ambiente=\env

	\newcounter{passo}	
	\newenvironment{passo}{%
		\refstepcounter{passo}%
		\penalty-200
		\noindent\colorbox{yellow!60!black}{\makebox[\columnwidth-\fboxsep*2][c]{\textbf{\color{white}Passo~\thepasso}}}\par\smallskip
	}{\par\medskip}

	% Papel de parede
	\CenterWallPaper{1}{../figuras/leao-pensador-wattermark.png}

	% Configuração extra: espaço entre as colunas.
	\setlength{\columnsep}{2cm}
	
	\renewcommand{\footnoterule}{{\null\hfill\rule{0.5\columnwidth}{0.5pt}\hfill\null}}
	
\begin{document}

	\begin{center}\includegraphics[width=0.8\columnwidth]{../figuras/LogotipoCursoLaTeX_v2}\end{center}

	\section*{Expressões matemáticas no \LaTeX\ --- Série 2}	
	\begin{center}
		\bfseries\color{red}
		Atenção: antes de fazer esta série, faça a primeira.
	\end{center}
		
	\begin{exercicio}
	Vamos complicar um pouco? Seu objetivo agora é construir a equação
	\[
	K_\nu (z) = \sqrt{\frac{\pi}{2z}} \frac{e^{-z}}{\left( \nu - \frac{1}{2} \right)!} \sum_{r = 0}^\infty
		\frac{ \left(\nu - \frac{1}{2}\right)! }{ r! \left( \nu - r - \frac{1}{2} \right)! } \left(2z\right)^{-r}
			\int\limits_0^\infty e^{-t}\, t^{\nu + r - 1/2}\, dt,
	\]
	mas você não fará isto sozinho, pelo menos não agora\rlap{.}\footnote{Você notou como a equação faz parte da oração? Além da vírgula após a equação, o parágrafo continua normalmente (a palavra ``mas'' não sofreu tabulação).} Depois de terminado, envie seu arquivo de instruções com a expressão acima para o professor através do moodle (você \emph{não} precisa enviar cada um dos passos abaixo).
	
	Independentemente do tamanho, toda expressão matemática é construída aos poucos. Então vamos juntos: siga as instruções dos passos abaixo e muito em breve você perceberá o macete.
	
	Mas antes, duas dicas: primeira, utilize \verb|\[| e \verb|\]| (estilo de exibição) ao invés de \verb|\(| e \verb|\)| (estilo de texto) para melhorar a aparência;\footnote{Você não está sob nenhum tipo de restrição vertical, então por que restringir sua expressão?} segunda, recompile seu arquivo a cada nova alteração e tenha certeza de ter acertado antes de continuar.
	
	Fique atento a dois erros que você fatalmente cometerá: primeiro, esquecer de abrir ou de fechar alguma chave. Segundo, errar o nome de um comando. Obserque que o TeXnicCenter coloca em azul os comandos que ele conhece. Isto serve para te ajudar. Portanto, use!
	
	Sem mais delongas, vamos lá! Crie um novo arquivo \texttt{tex} e\dots
	
	\bigskip
	
	%---------------------------------------------------------
	\begin{passo}\label{ex:A} Construa a seguinte expressão, que corresponde à extrema esquerda da equação acima:
	\[K_\nu(z) = \sqrt{\frac{\pi}{2z}}\]
	\end{passo}
	
	%---------------------------------------------------------
	\begin{passo}\label{ex:nu-1/2} Comece uma nova expressão (isto é, digite \verb|\[|\dots\verb|\]|) e reproduza o resultado abaixo.
	
	\[\left( \nu - \frac{1}{2} \right)!\]
	
	\noindent obs.\smallskip: não se esqueça de usar \verb|\left| e \verb|\right|.
	\end{passo}	
	
	%---------------------------------------------------------	
	\begin{passo}\label{ex:B}
	
	Comece outra nova expressão, mas desta vez, ao invés de digitá-la toda, copie o passo \ref{ex:nu-1/2} para o denominador (segundo argumento de \verb|\frac|):	
	\[
		\frac{e^{-z}}{\left( \nu - \frac{1}{2} \right)!}
	\]	
	\noindent obs.: observe como os parênteses se adequam automaticamente à nova situação.
	\end{passo}
	%---------------------------------------------------------
	
	\bigskip
	
	Perceba que a todo momento eu lhe peço para começar uma nova expressão. Deste modo você mantém a expressão simples e, por conseguinte, reduz a possibilidade de erros. Além disso, caso você cometa algum erro, é muito mais fácil encontrá-lo.
	
	\bigskip
	
	%---------------------------------------------------------
	\begin{passo}
	
	Comece esta nova expressão assim: \verb|\[\frac{}{}\]|. Em seguida, copie o passo \ref{ex:nu-1/2} para os \emph{dois} argumentos de \verb|\frac| e altere o segundo (denominador) para obter	
	\[
		\frac{ \left(\nu - \frac{1}{2}\right)! }{ r! \left( \nu - r - \frac{1}{2} \right)! }
	\]
	\end{passo}
	
	%---------------------------------------------------------		
	\begin{passo}
	
	Vamos trabalhar um pouco em cima da expressão que você criou no passo anterior.
	Simplesmente adicione $\sum_{r=0}^\infty$ à esquerda para obter	
	\[
	\sum_{r = 0}^\infty \frac{ \left(\nu - \frac{1}{2}\right)! }{ r! \left( \nu - r - \frac{1}{2} \right)! }
	\]
	\end{passo}	
	
	%---------------------------------------------------------	
	\begin{passo}\label{ex:C}
	
	À direita, coloque $(2z)^{-r}$:
	\[
		\sum_{r = 0}^\infty\frac{ \left(\nu - \frac{1}{2}\right)! }{ r! \left( \nu - r - \frac{1}{2} \right)! } (2z)^{-r}
	\]	
	\end{passo}
	%---------------------------------------------------------
	
	\bigskip
	
	Veja: esta é outra maneira de construir uma expressão. Você não precisa criar inúmeras expressões para depois copiá-las para outras\dots\ estamos fazendo isto simplesmente porque é mais fácil. Com um pouco de prática você passará a determinar as etapas de uma expressão mentalmente, mas enquanto isto não acontece, vamos por partes.
	
	\bigskip
	
	%---------------------------------------------------------
	\begin{passo}	
	Reproduza o resultado abaixo numa nova expressão (não se esqueça das chaves e observe que $1/2$ não usa \verb|\frac|)
	\[
	t^{\nu + r - 1/2}
	\]	
	\end{passo}	
		
	%---------------------------------------------------------
	\begin{passo} Na expressão do passo anterior, coloque $e^{-t}$ à esquerda:	
	\[
		e^{-t}\, t^{\nu + r - 1/2}
	\]
	dica: coloque o comando \verb|\,| (um pequeno espaço) entre $e^{-t}$ e $t^{\nu + \cdots}$.
	\end{passo}
	
	%---------------------------------------------------------	
	\begin{passo} Ainda à esquerda, use \verb|\int| para produzir	
	\[
		\int\limits_0^\infty e^{-t}\, t^{\nu + r - 1/2}
	\]	
	obs.: fica a seu critério utilizar \verb|\limits| ou não.
	\end{passo}
	
	%---------------------------------------------------------
	\begin{passo} Agora, adicione $dt$ à direita da expressão:	
	\[
	\int\limits_0^\infty e^{-t}\, t^{\nu + r - 1/2}\, dt
	\]	
	obs.: aqui é importante utilizar \verb|\,| antes de $dt$.
	\end{passo}	
	
	%---------------------------------------------------------
	\begin{passo} Agora, coloque o passo \ref{ex:C} à esquerda da integral:
	\[
		\sum_{r = 0}^\infty\frac{ \left(\nu - \frac{1}{2}\right)! }{ r! \left( \nu - r - \frac{1}{2} \right)! }
		(2z)^{-r} \int\limits_0^\infty e^{-t}\, t^{\nu + r - 1/2}\, dt
	\]
	dica: lembre-se de que o \LaTeX\ ignora os espaços que você digita. Use-os para organizar a expressão.
	\end{passo}
	
	%---------------------------------------------------------
	\begin{passo} Para terminar, copie os passos \ref{ex:A} e \ref{ex:B}, nesta ordem, para a esquerda:
	\[
		K_\nu(z) = \sqrt{\frac{\pi}{2z}} \frac{e^{-z}}{\left( \nu - \frac{1}{2} \right)!} \sum_{r = 0}^\infty\frac{
		\left(\nu - \frac{1}{2}\right)! }{ r! \left( \nu - r - \frac{1}{2} \right)! } (2z)^{-r}
		\int\limits_0^\infty e^{-t}\, t^{\nu + r - 1/2}\, dt
	\]
	\end{passo}
	
	\bigskip
	
	É isso aí! Não é tão difícil, é? Não se esqueça de enviar seu arquivo \texttt{tex}.
	\end{exercicio}
	
\end{document}